\documentclass[11pt,a4paper]{article}
\usepackage[utf8]{inputenc}
\usepackage[T1]{fontenc}
\usepackage{lmodern}
\usepackage[spanish]{babel}
\usepackage[margin=2.5cm]{geometry}
\usepackage{booktabs}
\usepackage{graphicx}
\usepackage{longtable}
\usepackage{float}
\usepackage[hidelinks]{hyperref}

\title{TA-IND-02: Priorización de requisitos con MoSCoW, Kano y votación ponderada\\
\large Dominio: Clínica veterinaria --- historias clínicas y vacunas}
\author{Jeanpierre Robinson Espinoza\\
Ingeniería de Requerimientos (ISR-401), UTEQ}
\date{Noviembre de 2026}

\begin{document}
\maketitle

\section{Dominio y requisitos}
El dominio trabajado es una \textbf{clínica veterinaria}, centrado en las historias clínicas y el control de vacunas. Es un dominio distinto al empleado en el PFC, por lo que no se reutilizan sus datos ni sus instrumentos. El autor formuló veinte requisitos: doce funcionales y ocho de calidad. Cada requisito incluye un criterio verificable con métrica, umbral y condición de medición, siguiendo las recomendaciones de \cite{wiegers2013} y \cite{iso29148}. El esfuerzo total tabulado es de 114 horas, sin que se modificaran las horas originales para ajustar categorías.

\textbf{Requisitos formulados: 20 en total, 12 funcionales y 8 de calidad.}\\

\begin{longtable}{@{}llp{10.5cm}r@{}}
\toprule
ID & Tipo & Enunciado & h \\
\midrule
\endhead
R01 & F & El sistema deberá permitir registrar cada mascota con un código único y los campos obligatorios nombre, especie y propietario; además podrá guardar raza, sexo, edad, peso y señas particulares. & 8 \\
R02 & F & El sistema deberá permitir registrar los datos de contacto del propietario, al menos nombre y teléfono, asociados a una o varias mascotas. & 4 \\
R03 & F & El sistema deberá permitir registrar antecedentes, alergias y enfermedades previas de cada mascota. & 5 \\
R04 & F & El sistema deberá permitir registrar cada consulta con los seis campos definidos: fecha, motivo, síntomas, diagnóstico, tratamiento y observaciones del veterinario. & 8 \\
R05 & F & El sistema deberá mostrar en orden cronológico todas las consultas anteriores de una mascota. & 6 \\
R06 & F & El sistema deberá permitir registrar las vacunas aplicadas a cada mascota, indicando tipo de vacuna, fecha de aplicación y responsable. & 6 \\
R07 & F & El sistema deberá permitir registrar la fecha de la próxima dosis o del siguiente control de vacunación. & 4 \\
R08 & F & El sistema deberá mostrar al personal de la clínica un aviso visible al menos 3 días antes de la fecha programada de vacunación, cuando exista una próxima dosis registrada. & 6 \\
R09 & F & El sistema deberá mostrar una lista de mascotas cuya fecha de vacunación sea igual o anterior a los 7 días siguientes, para planificar los controles de la semana. & 5 \\
R10 & F & El sistema deberá permitir buscar mascotas por nombre, código de paciente o teléfono del propietario. & 5 \\
R11 & F & El sistema deberá gestionar citas veterinarias registrando fecha, hora, mascota y motivo de atención en un único registro consultable. & 6 \\
R12 & F & El sistema deberá registrar el seguimiento complementario de cada mascota (desparasitación y control de peso) como eventos fechados con observación. & 4 \\
R13 & C & El sistema deberá restringir la creación, modificación y eliminación de historias clínicas según el rol del usuario. & 8 \\
R14 & C & El sistema deberá garantizar la trazabilidad de las modificaciones de historias clínicas, registrando el usuario responsable y la fecha y hora de cada cambio. & 6 \\
R15 & C & El sistema deberá mantener la integridad de los registros evitando duplicados por código de mascota o por coincidencia de nombre, especie y teléfono del propietario. & 6 \\
R16 & C & El sistema deberá permitir consultar una ficha clínica en máximo 3 segundos en el 95 \% de los casos, con hasta 1.000 mascotas registradas y conexión estable. & 5 \\
R17 & C & El sistema deberá realizar al menos una copia de seguridad diaria de la información clínica y permitir recuperar los datos respaldados en máximo 4 horas. & 6 \\
R18 & C & El sistema deberá permitir que al menos 4 de cada 5 usuarios del personal completen el registro y la búsqueda de una mascota en máximo 3 minutos, sin asistencia, en una prueba de uso. & 6 \\
R19 & C & El sistema deberá mantener la portabilidad del historial clínico, permitiendo exportarlo o imprimirlo con datos completos y legibles para entregarlo al propietario o compartirlo con otro profesional. & 4 \\
R20 & C & El sistema deberá mantener disponibles los registros durante el horario de atención definido y configurable por la clínica, con al menos 99 \% mensual. & 6 \\
\bottomrule
\end{longtable}



Las entrevistas y el cuestionario Kano se aplicaron con los códigos \texttt{P1}--\texttt{P5}. Las justificaciones citan esos códigos solo cuando el original entrevistado respalda la afirmación.

\section{Priorización con MoSCoW}
La clasificación MoSCoW organiza el esfuerzo en \textbf{Must} (imprescindible para la primera versión), \texttt{Should} (importante pero no bloqueante), \texttt{Could} (deseable si sobra capacidad) y \texttt{Won't} (fuera de la versión 1). Se aplicó la regla de que los Must no superen el 60\% del esfuerzo total. El límite para este conjunto es 68,4 h; la versión corregida queda en 66 h (57,89\%).

\begin{table}[H]
\centering
\footnotesize
\caption{Distribución del esfuerzo por categoría MoSCoW (114 h en total).}
\label{tab:moscow}
\begin{tabular}{@{}lrrr@{}}
\toprule
Categoría & Requisitos & Horas & Porcentaje \\
\midrule
Could (C) & 1 & 4 & 3,51 \% \\
Must (M) & 11 & 66 & 57,89 \% \\
Should (S) & 7 & 40 & 35,09 \% \\
Won't (W) & 1 & 4 & 3,51 \% \\
\bottomrule
\end{tabular}
\end{table}


Para cumplir la regla sin rebajar horas se reclasificaron dos requisitos: \texttt{R09} pasó de Must a Should porque el control de vacunas pendientes quedó empatado entre atractivo e indiferente en Kano; \texttt{R12} pasó a Won't por ser un seguimiento complementario y no imprescindible. El total de horas se conserva en 114. El requisito \texttt{R19} se formuló como atributo de portabilidad, no como una función aislada, de modo que se mantienen ocho requisitos de calidad.

\section{Análisis Kano}
Se aplicaron ocho características Kano a cinco participantes (\texttt{P1}--\texttt{P5}). Cada característica tiene una pregunta funcional y otra disfuncional, lo que produce 40 pares de respuestas y 80 respuestas en total. La clasificación usa la tabla de evaluación estándar: atractiva (A), indiferente (I), basicidad (M), desempeño (O), reversa (R) y cuestionable (Q). El enfoque sigue la lógica de priorización descrita por \cite{pohl2025}.

\begin{table}[H]
\centering
\footnotesize
\caption{Correspondencia entre las ocho características Kano y los requisitos enlazados.}
\label{tab:kano}
\begin{tabular}{@{}lll@{}}
\toprule
ID & Característica & Requisito \\
\midrule
K1 & Registro de mascotas & R01 \\
K2 & Historial de consultas & R05 \\
K3 & Registro de vacunas & R06 \\
K4 & Avisos de próximas vacunas & R08 \\
K5 & Búsqueda de historias clínicas & R10 \\
K6 & Control de vacunas pendientes & R09 \\
K7 & Acceso restringido & R13 \\
K8 & Descarga del historial & R19 \\
\bottomrule
\end{tabular}
\end{table}


\textbf{Clasificaci\'{o}n Kano (5 participantes). A=atractiva, I=indiferente, M=basicidad, O=desempe\~{n}o, R=reversa, Q=cuestionable.}\\

\begin{tabular}{@{}lrrrrrrl@{}}
\toprule
Caracter\'{i}stica & A & I & M & O & R & Q & Resultado \\
\midrule
K1 & 1 & 1 & 0 & 3 & 0 & 0 & O \\
K2 & 2 & 0 & 0 & 3 & 0 & 0 & O \\
K3 & 1 & 0 & 0 & 4 & 0 & 0 & O \\
K4 & 2 & 3 & 0 & 0 & 0 & 0 & I \\
K5 & 2 & 0 & 0 & 3 & 0 & 0 & O \\
K6 & 2 & 2 & 0 & 1 & 0 & 0 & empate I/A \\
K7 & 0 & 1 & 0 & 4 & 0 & 0 & O \\
K8 & 0 & 5 & 0 & 0 & 0 & 0 & I \\
\bottomrule
\end{tabular}



Las características con mayoría desempeño son \texttt{K1}, \texttt{K2}, \texttt{K3}, \texttt{K5} y \texttt{K7}; \texttt{K4} y \texttt{K8} son indiferentes. \texttt{K6} no tiene mayoría: dos respuestas atractivas, dos indiferentes y una de desempeño. Por eso se conserva el empate y no se fuerza una categoría; ese resultado es la base para mantener \texttt{R09} como Should.

\section{Votación ponderada}
Se recibieron las tres boletas \texttt{V1}--\texttt{V3}, aplicadas sobre los veinte enunciados finales. La lectura directa mostró 101 puntos en \texttt{V1} y en \texttt{V3}, aunque las tres boletas imprimen «TOTAL 100». El autor confirmó dos correcciones: \texttt{V1} bajó \texttt{R20} de 3 a 2 puntos y \texttt{V3} bajó \texttt{R16} de 4 a 3 puntos. No se modificó ningún otro valor. Después de esas correcciones, cada boleta suma exactamente 100 y \texttt{datos/votacion.csv} contiene 60 filas (3 votantes $\times$ 20 requisitos).

\begin{table}[H]
\centering
\footnotesize
\caption{Votación ponderada: totales por requisito y rangos promediados.}
\label{tab:votacion}
\begin{tabular}{@{}lrrrrr@{}}
\toprule
ID & V1 & V2 & V3 & Total & Rango \\
\midrule
R01 & 9 & 6 & 5 & 20 & 18,5 \\
R02 & 4 & 5 & 6 & 15 & 12,0 \\
R03 & 4 & 4 & 6 & 14 & 8,0 \\
R04 & 8 & 7 & 5 & 20 & 18,5 \\
R05 & 4 & 5 & 4 & 13 & 5,0 \\
R06 & 8 & 7 & 6 & 21 & 20,0 \\
R07 & 6 & 6 & 5 & 17 & 16,5 \\
R08 & 4 & 5 & 6 & 15 & 12,0 \\
R09 & 4 & 4 & 5 & 13 & 5,0 \\
R10 & 4 & 6 & 5 & 15 & 12,0 \\
R11 & 6 & 5 & 6 & 17 & 16,5 \\
R12 & 4 & 4 & 5 & 13 & 5,0 \\
R13 & 5 & 5 & 4 & 14 & 8,0 \\
R14 & 5 & 6 & 5 & 16 & 15,0 \\
R15 & 5 & 4 & 5 & 14 & 8,0 \\
R16 & 4 & 4 & 3 & 11 & 2,0 \\
R17 & 5 & 5 & 5 & 15 & 12,0 \\
R18 & 4 & 4 & 4 & 12 & 3,0 \\
R19 & 5 & 5 & 5 & 15 & 12,0 \\
R20 & 2 & 3 & 5 & 10 & 1,0 \\
\bottomrule
\end{tabular}
\end{table}


El rho de Spearman entre el orden MoSCoW y el total de votos, sobre los veinte requisitos y con rangos promediados, es $\rho = 0{,}3904$. El cálculo aplica Pearson sobre los rangos y valida que cada votante sume 100. Los requisitos con mayor puntaje son \texttt{R06} (21 puntos), \texttt{R01} y \texttt{R04} (20 puntos); los de menor puntaje son \texttt{R20} (10), \texttt{R16} (11) y \texttt{R18} (12).

\section{Comparación de técnicas}
MoSCoW aporta una restricción de capacidad: 114 h en total y un techo de 66 h para Must. Kano aporta la percepción del usuario: cinco de ocho características muestran desempeño, dos son indiferentes y una queda empatada. La votación ponderada aporta el peso relativo que los tres votantes asignan a los mismos veinte requisitos. Las tres técnicas no se sustituyen entre sí: MoSCoW decide factibilidad, Kano explica satisfacción y la votación ponderada mide prioridad percibida.

\textbf{Comparaci\'{o}n de los tres \'{o}rdenes sobre los 8 requisitos enlazados con Kano.}\\

\begin{tabular}{@{}lrrrrrr@{}}
\toprule
ID & M & Kano & Voto & Rango M & Rango K & Rango V \\
\midrule
R01 & 4 & 0.80 & 20 & 6.0 & 4.5 & 7.0 \\
R05 & 4 & 1.00 & 13 & 6.0 & 7.0 & 1.5 \\
R06 & 4 & 1.00 & 21 & 6.0 & 7.0 & 8.0 \\
R08 & 3 & 0.40 & 15 & 2.5 & 2.0 & 5.0 \\
R09 & 3 & 0.60 & 13 & 2.5 & 3.0 & 1.5 \\
R10 & 4 & 1.00 & 15 & 6.0 & 7.0 & 5.0 \\
R13 & 4 & 0.80 & 14 & 6.0 & 4.5 & 3.0 \\
R19 & 2 & 0.00 & 15 & 1.0 & 1.0 & 5.0 \\
\bottomrule
\end{tabular}



Para comparar los tres órdenes sobre el mismo conjunto se usaron los ocho requisitos enlazados con Kano. El orden Kano se derivó del coeficiente \emph{better} $(A+O)/5$, que conserva el empate de \texttt{K6} sin forzar una categoría; MoSCoW usa \texttt{M}=4, \texttt{S}=3, \texttt{C}=2 y \texttt{W}=1, y la votación usa el total de puntos. Los rho de Spearman con rangos promediados son: MoSCoW--Kano $0{,}8930$, MoSCoW--voto $0{,}1772$ y Kano--voto $0{,}1013$. MoSCoW y Kano coinciden ampliamente, pero la votación ponderada se incrementa de forma distinta: los votantes prefieren \texttt{R06} y \texttt{R01}, mientras que Kano valora también \texttt{R05}, \texttt{R10} y \texttt{R13}. Estas diferencias se reportan sin forzar acuerdos artificiales.

\section{Decisión final}
La versión 1 mantiene diecinueve requisitos: los once Must (66 h), los siete Should (40 h) y el único Could (\texttt{R19}, 4 h). Queda excluido \texttt{R12} (W, 4 h) porque registra desparasitación y control de peso como seguimiento complementario. La votación no revierte esa exclusión: \texttt{R12} recibió 13 puntos, por debajo de \texttt{R06} (21) y \texttt{R01} (20), y comparte el extremo bajo del orden con \texttt{R05} y \texttt{R09}.

\texttt{R09} sigue como Should: el empate de \texttt{K6} y la existencia de \texttt{R08} con aviso individual de 3 días justifican el horizonte semanal de 7 días. La decisión final combina la restricción de esfuerzo, la satisfacción medida por Kano y la prioridad percibida en la votación. Ninguna herramienta de IA participó en esta decisión ni en la elección de exclusiones.

\section{Declaración de uso de inteligencia artificial}
El autor, Jeanpierre Robinson Espinoza, definió el dominio, redactó los veinte requisitos, sus criterios de verificación y sus justificaciones, y es responsable de aplicar los instrumentos a las personas participantes (\texttt{P1}--\texttt{P5} en Kano y \texttt{V1}--\texttt{V3} en votación), de la tabulación, de comprobar los cálculos y de la decisión final. Confirmó que las cinco respuestas Kano son de personas reales y confirmó las dos correcciones de una casilla en las boletas \texttt{V1} y \texttt{V3} después de detectar que sumaban 101 puntos. OpenCode (modelo Big Pickle) y Claude (Anthropic) actuaron únicamente como apoyo técnico bajo su supervisión: transcribir los instrumentos a LaTeX, preparar plantillas y boletas en blanco, adaptar el esquema de los CSV, convertir \texttt{Must}/\texttt{Should}/\texttt{Could} a \texttt{M}/\texttt{S}/\texttt{C}, escribir los scripts de clasificación Kano y de Spearman con rangos promediados, corregir errores de compilación y revisar la redacción. La IA no generó participantes, respuestas, votos, requisitos ni decisiones; las justificaciones citan \texttt{P1}--\texttt{P5} solo donde constan en los originales. Esta revisión también se declara, y el historial de commits conserva íntegramente los cambios realizados.

\bibliographystyle{IEEEtran}
\bibliography{referencias}

\end{document}
